% problem-000425 9 有机碱性比较
\section*{9. 有机碱性比较}
\textit{下列为分问。}

\subsection*{9-1}
下列化合物中碱性最强的是

（图：）

A

（图：）

B

（图：）

C

（图：）

D

\subsection*{9-2}
下列化合物在酸性条件下最难发⽣⽔解的是

（图：）

A

（图：）

B

（图：）

C

（图：）

D

\subsection*{9-3}
部分丙烷与⻧素⾃由基发⽣的反应其相应的反应焓变如下：$\mathrm { C H _ { 3 } C H _ { 2 } C H _ { 3 } + C l \cdot  C H _ { 3 } C H C H _ { 3 } + H C l }$ $\Delta H ^ { \ominus } = - 3 4 \mathrm { k J } \mathrm { m o l ^ { - 1 } }$ (1)$\mathrm { C H _ { 3 } C H _ { 2 } C H _ { 3 } + C l \cdot  C H _ { 3 } C H _ { 2 } C H _ { 2 } + H C l }$ $\Delta { H ^ { \ominus } } \ = \ - 2 2 \mathrm { k J } \mathrm { m o l ^ { - 1 } }$ (2)$\mathrm { C H _ { 3 } C H _ { 2 } C H _ { 3 } } + \mathrm { B r } \cdot  \mathrm { C H _ { 3 } C H _ { 2 } C H _ { 2 } } + \mathrm { H B r }$ $\Delta H ^ { \ominus } = + 4 4 \mathrm { k J m o l ^ { - 1 } }$ (3)

当反应温度升⾼，上述反应的速率将A．都降低 B．都升⾼ C．反应(1)和(2)降低，反应(3)升⾼ D．反应(1)和(2)升⾼，反应(3)降低

\subsection*{9-4}
$2 5 ~ ^ { \circ } \mathrm { C } .$ ，下列反应两种产物的⽐例已经给出43% 57%$\mathrm { C H _ { 3 } C H _ { 2 } C H _ { 3 } + C l \cdot  C H _ { 3 } C H C H _ { 3 } + C H _ { 3 } C H _ { 2 } C H _ { 2 } + H C l }$

结合9-3题中给出的信息，预测当反应温度升⾼⽽其他反应条件不变时，1级⾃由基产物的⽐例将A．降低 B．升⾼ C．不变 D．⽆法判断
